%%% BEGIN fig:ee-cutsum-disc
\begin{figure}[htbp]
\centering
% three-loop ladder graph of the current correlator (origin at its centre); #1 = cut path
\newcommand{\fbEEladder}[1]{%
  \draw[photon] (-1.85,0) -- (-1.25,0);
  \draw[photon] (1.25,0) -- (1.85,0);
  \draw[plain] (-1.25,0) to[out=62,in=180] (-0.62,0.72) -- (-0.42,0.72);
  \draw[plain] (-0.42,0.72) -- (0.42,0.72);
  \draw[quark=0.55] (0.42,0.72) -- (0.62,0.72) to[out=0,in=118] (1.25,0);
  \draw[plain] (1.25,0) to[out=-118,in=0] (0.62,-0.72) -- (0.42,-0.72);
  \draw[plain] (0.42,-0.72) -- (-0.42,-0.72);
  \draw[quark=0.55] (-0.42,-0.72) -- (-0.62,-0.72) to[out=180,in=-62] (-1.25,0);
  \draw[gluon] (-0.42,0.72) -- (-0.42,-0.72);
  \draw[gluon] (0.42,0.72) -- (0.42,-0.72);
  \foreach \p in {(-1.25,0),(1.25,0),(-0.42,0.72),(0.42,0.72),(-0.42,-0.72),(0.42,-0.72)} \node[vertex] at \p {};
  \draw[cut] #1;}
% reduced diagram of e+e- (origin at its centre); #1 = 1 draws the cut, 0 the forward box
\newcommand{\fbEEred}[1]{%
  \node[hard blob,circle,minimum size=0.42cm,inner sep=0pt] (H)  at (-1.13,0) {};
  \node[hard blob,circle,minimum size=0.42cm,inner sep=0pt] (Hd) at ( 1.13,0) {};
  \node[jet blob,ellipse,minimum width=1.42cm,minimum height=0.46cm,inner sep=0pt] (J1) at (0, 0.78) {};
  \node[jetB blob,ellipse,minimum width=1.42cm,minimum height=0.46cm,inner sep=0pt] (J2) at (0,-0.78) {};
  \node[soft blob,circle,minimum size=0.46cm,inner sep=0pt] (S) at (0,0) {};
  \draw[photon] (-1.5,0) -- (H.west);
  \draw[photon] (Hd.east) -- (1.5,0);
  \draw[quark] (H.70) -- (J1.188);
  \draw[quark] (J1.352) -- (Hd.110);
  \draw[quark] (Hd.250) -- (J2.8);
  \draw[quark] (J2.172) -- (H.290);
  \draw[gluon] (S.140) -- (-0.36,0.56);
  \draw[gluon] (S.40) -- (0.36,0.56);
  \draw[gluon] (S.220) -- (-0.36,-0.56);
  \draw[gluon] (S.320) -- (0.36,-0.56);
  \ifnum#1=1 \node[font=\scriptsize,inner sep=0pt,anchor=east] at ($(S.west)+(-0.04,0)$) {$S$};\else \node[font=\scriptsize] at (S) {$S$};\fi
  \node[font=\scriptsize] at (-0.47,0.78) {$J_1$};
  \node[font=\scriptsize] at (-0.47,-0.78) {$J_2$};
  \node[font=\scriptsize,anchor=north] at (H.south) {$H$};
  \node[font=\scriptsize,anchor=north] at (Hd.south) {$H^\dagger$};
  \ifnum#1=1 \draw[cut] (0,1.18) -- (0,-1.18);
  \else \draw[draw=gray!60!black,line width=0.6pt,dash pattern=on 2.5pt off 1.5pt,rounded corners=4pt]
          (-0.9,-1.14) rectangle (0.9,1.14);
  \fi}
\begin{tikzpicture}[fd,
  fb lab/.style={font=\footnotesize},
  fb note/.style={font=\scriptsize,text=red!55!black},
  fb bcut/.style={draw,line width=0.9pt,decorate,decoration={zigzag,segment length=4.2pt,amplitude=1.6pt}},
  fb contour/.style={draw=blue!60!black,line width=0.9pt},
  fb carrow/.style={postaction={decorate,decoration={markings,
      mark=at position #1 with {\arrow{Stealth[length=5.5pt,width=4.5pt]}}}}}]
  % ======================= (a) cut classes of one O(alpha_s^2) graph =======================
  \begin{scope}
    \fbEEladder{(0,1.05) -- (0,-1.05)}
    \node[fb lab] at (0,-1.38) {两粒子割线 $q\bar q$};
    \node[fb note] at (0,-1.76) {单类至多 $1/\varepsilon^4$ 极点};
  \end{scope}
  \begin{scope}[xshift=4.35cm]
    \fbEEladder{(-0.79,0.96) -- (0.09,-0.96)}
    \node[fb lab] at (0,-1.38) {三粒子割线 $q\bar qg$};
    \node[fb note] at (0,-1.76) {单类至多 $1/\varepsilon^4$ 极点};
  \end{scope}
  \begin{scope}[xshift=8.7cm]
    \fbEEladder{(-0.91,0.965) -- (0.91,-0.965)}
    \node[fb lab] at (0,-1.38) {四粒子割线 $q\bar qgg$};
    \node[fb note] at (0,-1.76) {单类至多 $1/\varepsilon^4$ 极点};
  \end{scope}
  % hand-made brace and the statement
  \draw[gray!45!black,line width=0.7pt] (-1.7,-2.0) to[out=-90,in=180] (-1.5,-2.2) -- (4.15,-2.2)
       to[out=0,in=90] (4.35,-2.4) to[out=90,in=180] (4.55,-2.2) -- (10.2,-2.2) to[out=0,in=-90] (10.4,-2.0);
  \node[font=\small] at (4.35,-2.75) {对全部割线求和（连同 $\MSbar$ 抵消项）：红外极点相消 $\;\Longrightarrow\;$ 有限};
  \node[font=\footnotesize,text=gray!30!black] at (4.35,-3.25) {（单态项 $\propto(\sum_qe_q)^2$ 自 $\order{\as^3}$ 起：Furry 定理）};
  \node[font=\small,anchor=west] at (-2.0,1.2) {(a)};
  % ======================= (b) cut sum = discontinuity of the forward function =======================
  \begin{scope}[yshift=-5.7cm]
    \node[font=\small,anchor=west] at (-2.0,1.45) {(b)};
    \node[font=\large] at (-1.55,0) {$\displaystyle\sum_{C}$};
    \begin{scope}[xshift=0.4cm]\fbEEred{1}\end{scope}
    \node[font=\normalsize] at (2.78,0) {$=\ \Disc_s$};
    \begin{scope}[xshift=5.25cm]\fbEEred{0}\end{scope}
    \node[fb lab,anchor=north] at (0.4,-1.28) {割线只穿过 $J_1$、$J_2$、$S$};
    \node[fb lab,anchor=north,align=center] at (5.25,-1.28) {前向（编时）组合 $F(s)$};
    % complex s plane
    \begin{scope}[xshift=10.0cm]
      \draw[lcaxis] (-1.95,0) -- (2.7,0) node[right,font=\footnotesize,inner sep=2pt] {$\Real s$};
      \draw[lcaxis] (0,-1.8) -- (0,1.95) node[left,font=\footnotesize] {$\Imag s$};
      \draw[fb bcut] (0,0) -- (2.45,0);
      \draw[fb contour,fb carrow=0.2,fb carrow=0.5,fb carrow=0.8] (5:1.45) arc[start angle=5,end angle=355,radius=1.45];
      \fill[blue!60!black] (5:1.45) circle[radius=1.5pt];
      \fill[blue!60!black] (-5:1.45) circle[radius=1.5pt];
      \node[font=\scriptsize,anchor=south west,inner sep=1pt] at ($(5:1.45)+(0.04,0.07)$) {$s+\ii0$};
      \node[font=\scriptsize,anchor=north west,inner sep=1pt] at ($(-5:1.45)+(0.04,-0.07)$) {$s-\ii0$};
      \node[font=\scriptsize,text=blue!55!black,anchor=south east,inner sep=1pt] at (132:1.5) {$|z|=s$：无钳夹};
      \node[font=\scriptsize,align=left,text=gray!25!black,anchor=north west,inner sep=1pt] at (1.62,-0.62) {钳夹（Landau\\奇点）只在\\实轴上};
      \node[font=\scriptsize,anchor=north] at (0.35,-1.95) {$F(s+\ii0)-F(s-\ii0)=-\displaystyle\int_{|z|=s}\!\dd z\,F'(z)$};
    \end{scope}
  \end{scope}
\end{tikzpicture}
\caption{割线求和与前向函数的不连续性（$e^+e^-$ 总截面；红色虚线为末态割线，左侧为振幅，右侧为共轭振幅）。
(a)~$\order{\as^2}$ 的一个三圈流关联函数图（夸克圈上两条胶子横档）的三类割线：两粒子 $q\bar q$、三粒子 $q\bar qg$ 与四粒子 $q\bar qgg$（每类只画一例，其余割线位置与镜像未画出）。
每类割图单独含有直到 $1/\varepsilon^4$ 的软、共线极点，只有对全部割线求和（连同 $\MSbar$ 抵消项）才有限（\cref{cor:ee-cuts}）；
两个夸克圈各连一个光子的单态项 $\propto(\sum_qe_q)^2$ 由 Furry 定理自 $\order{\as^3}$ 起才出现。
(b)~在钳夹面附近，割线只能穿过末态喷注 $J_1$、$J_2$ 与软子图 $S$，不能穿过 $H$、$H^\dagger$；对固定的 $H$、$H^\dagger$，由幺正性，对这些割线求和等于不加割线的前向（编时）组合 $F$（虚线框）的不连续性 $F(s+\ii0)-F(s-\ii0)$。
$F$ 是 $s$ 的解析函数，钳夹（Landau 奇点）只出现在实轴的割线上；由 Cauchy 定理，不连续性等于沿圆周 $|z|=s$ 的积分，圆周除两端外离实轴 $\order{s}$，那里没有钳夹，只剩欧氏区可计算的部分（围道公式，\cref{prop:ee-contour}）。}
\label{fig:ee-cutsum-disc}
\end{figure}
%%% END fig:ee-cutsum-disc

%%% BEGIN fig:ee-oneloop
\begin{figure}[htbp]
\centering
\begin{tikzpicture}[fd,
  t2 lab/.style={font=\small},
  t2 cutlab/.style={font=\scriptsize,text=red!65!black},
  t2 cut0/.style={cut,draw=red!30},
  t2 cutlab0/.style={font=\scriptsize,text=red!35}]
  % ======================= row 1: amplitudes =======================
  % (a) virtual vertex correction
  \begin{scope}
    \coordinate (V) at (0,0); \coordinate (E1) at (2.0,1.15); \coordinate (E2) at (2.0,-1.15);
    \coordinate (G1) at ($(V)!0.56!(E1)$); \coordinate (G2) at ($(V)!0.56!(E2)$);
    \draw[photon] (-1.4,0) -- (V);
    \draw[quark] (V) -- (G1); \draw[quark] (G1) -- (E1);
    \draw[quark] (E2) -- (G2); \draw[quark] (G2) -- (V);
    \draw[gluon] (G1) -- (G2);
    \foreach \p in {V,G1,G2} \node[vertex] at (\p) {};
    \path[mom={$q$},mom sep=0.28cm] (-1.4,0) -- (V);
    \path[mom={$p_1$},mom pos=0.6] (G1) -- (E1);
    \path[mom'={$p_2$},mom pos=0.6] (G2) -- (E2);
    \node at (0.3,-1.75) {(a)\ $\mathcal M_V$};
  \end{scope}
  % (b) real emission off the quark
  \begin{scope}[xshift=5.0cm]
    \coordinate (V) at (0,0); \coordinate (E1) at (2.0,1.15); \coordinate (E2) at (2.0,-1.15);
    \coordinate (G) at ($(V)!0.45!(E1)$);
    \draw[photon] (-1.4,0) -- (V);
    \draw[quark] (V) -- (G); \draw[quark] (G) -- (E1);
    \draw[quark] (E2) -- (V);
    \draw[gluon] (G) -- (2.1,0.05);
    \foreach \p in {V,G} \node[vertex] at (\p) {};
    \path[mom={$q$},mom sep=0.28cm] (-1.4,0) -- (V);
    \path[mom={$p_1$},mom pos=0.55] (G) -- (E1);
    \path[mom'={$p_2$},mom pos=0.7] (V) -- (E2);
    \path[mom'={$k$},mom pos=0.6] (G) -- (2.1,0.05);
    \node at (0.3,-1.75) {(b)\ $\mathcal M_q$};
  \end{scope}
  % (c) real emission off the antiquark
  \begin{scope}[xshift=10.0cm]
    \coordinate (V) at (0,0); \coordinate (E1) at (2.0,1.15); \coordinate (E2) at (2.0,-1.15);
    \coordinate (G) at ($(V)!0.45!(E2)$);
    \draw[photon] (-1.4,0) -- (V);
    \draw[quark] (V) -- (E1);
    \draw[quark] (E2) -- (G); \draw[quark] (G) -- (V);
    \draw[gluon] (G) -- (2.1,-0.05);
    \foreach \p in {V,G} \node[vertex] at (\p) {};
    \path[mom={$q$},mom sep=0.28cm] (-1.4,0) -- (V);
    \path[mom={$p_1$},mom pos=0.7] (V) -- (E1);
    \path[mom'={$p_2$},mom pos=0.55] (G) -- (E2);
    \path[mom={$k$},mom pos=0.6] (G) -- (2.1,-0.05);
    \node at (0.3,-1.75) {(c)\ $\mathcal M_{\bar q}$};
  \end{scope}
  % ======================= row 2: two-loop vacuum polarization with cuts =======================
  % loop: centre C=(0.3,0), radii 1.15 x 0.85; L=(-0.85,0), R=(1.45,0); fermion number flows L -> top -> R -> bottom -> L
  \begin{scope}[yshift=-4.3cm]
    % (d) vertex type
    \coordinate (L) at (-0.85,0); \coordinate (R) at (1.45,0);
    \coordinate (T) at ($(0.3,0)+(125:1.15 and 0.85)$); \coordinate (B) at ($(0.3,0)+(-55:1.15 and 0.85)$);
    \draw[photon] (-1.75,0) -- (L); \draw[photon] (R) -- (2.35,0);
    \draw[quark=0.7] (L) arc[start angle=180,end angle=0,x radius=1.15,y radius=0.85];
    \draw[quark=0.7] (R) arc[start angle=0,end angle=-180,x radius=1.15,y radius=0.85];
    \draw[gluon] (T) -- (B);
    \foreach \p in {L,R,T,B} \node[vertex] at (\p) {};
    \path[mom={$q$},mom sep=0.28cm] (-1.75,0) -- (L);
    \draw[cut] (-0.62,-1.6) -- (-0.62,1.6) node[above,t2 cutlab] {2};
    \draw[cut] (0.3,-1.6) -- (0.3,1.6) node[above,t2 cutlab] {3};
    \draw[cut] (1.22,-1.6) -- (1.22,1.6) node[above,t2 cutlab] {2};
    \node at (0.3,-2.15) {(d)};
  \end{scope}
  \begin{scope}[xshift=5.0cm,yshift=-4.3cm]
    % (e) self-energy on the quark line
    \coordinate (L) at (-0.85,0); \coordinate (R) at (1.45,0);
    \coordinate (T1) at ($(0.3,0)+(125:1.15 and 0.85)$); \coordinate (T2) at ($(0.3,0)+(55:1.15 and 0.85)$);
    \draw[photon] (-1.75,0) -- (L); \draw[photon] (R) -- (2.35,0);
    \draw[quark=0.61] (L) arc[start angle=180,end angle=0,x radius=1.15,y radius=0.85];
    \draw[quark=0.67] (R) arc[start angle=0,end angle=-180,x radius=1.15,y radius=0.85];
    \draw[gluon] (T1) .. controls ($(T1)+(0.05,0.75)$) and ($(T2)+(-0.05,0.75)$) .. (T2);
    \foreach \p in {L,R,T1,T2} \node[vertex] at (\p) {};
    \path[mom={$q$},mom sep=0.28cm] (-1.75,0) -- (L);
    \draw[t2 cut0] (-0.62,-1.6) -- (-0.62,1.6) node[above,t2 cutlab0] {2};
    \draw[cut] (0.3,-1.6) -- (0.3,1.6) node[above,t2 cutlab] {3};
    \draw[t2 cut0] (1.22,-1.6) -- (1.22,1.6) node[above,t2 cutlab0] {2};
    \node at (0.3,-2.15) {(e)};
  \end{scope}
  \begin{scope}[xshift=10.0cm,yshift=-4.3cm]
    % (f) self-energy on the antiquark line
    \coordinate (L) at (-0.85,0); \coordinate (R) at (1.45,0);
    \coordinate (B1) at ($(0.3,0)+(-125:1.15 and 0.85)$); \coordinate (B2) at ($(0.3,0)+(-55:1.15 and 0.85)$);
    \draw[photon] (-1.75,0) -- (L); \draw[photon] (R) -- (2.35,0);
    \draw[quark=0.33] (L) arc[start angle=180,end angle=0,x radius=1.15,y radius=0.85];
    \draw[quark=0.39] (R) arc[start angle=0,end angle=-180,x radius=1.15,y radius=0.85];
    \draw[gluon] (B2) .. controls ($(B2)+(-0.05,-0.75)$) and ($(B1)+(0.05,-0.75)$) .. (B1);
    \foreach \p in {L,R,B1,B2} \node[vertex] at (\p) {};
    \path[mom={$q$},mom sep=0.28cm] (-1.75,0) -- (L);
    \draw[t2 cut0] (-0.62,-1.6) -- (-0.62,1.6) node[above,t2 cutlab0] {2};
    \draw[cut] (0.3,-1.6) -- (0.3,1.6) node[above,t2 cutlab] {3};
    \draw[t2 cut0] (1.22,-1.6) -- (1.22,1.6) node[above,t2 cutlab0] {2};
    \node at (0.3,-2.15) {(f)};
  \end{scope}
\end{tikzpicture}
\caption{$e^+e^-\to$ 强子总截面的 $\order{\as}$ 修正。上排：振幅层次的图——(a)~单圈顶点修正 $\mathcal M_V$，(b)、(c)~夸克或反夸克辐射一个实胶子 $\mathcal M_q$、$\mathcal M_{\bar q}$
（无质量外线的自能修正在维数正规化中是无标度积分，为零，故未画出）。下排：电磁流关联函数 $\Pi$ 的三个两圈图及其全部切割（虚线上方的数字为被切粒子数）。
(d)~两粒子切割给出 $2\Real(\mathcal M_0^*\mathcal M_V)$（$\mathcal M_0$ 为 Born 振幅），三粒子切割给出干涉项 $2\Real(\mathcal M_q\mathcal M_{\bar q}^*)$；
(e)、(f)~三粒子切割给出 $|\mathcal M_q|^2$、$|\mathcal M_{\bar q}|^2$，两粒子切割（浅色）对应在壳无质量外线上的自能，在维数正规化中为零。
由切割规则，所有切割之和等于 $2\Imag\Pi$ 在 $\order{\as}$ 的贡献：虚、实两部分各含 $1/\varepsilon^2$ 极点，其和有限，给出 $R=R_0\big(1+\tfrac34C_F\,\as/\pi\big)$。}
\label{fig:ee-oneloop}
\end{figure}
%%% END fig:ee-oneloop

%%% BEGIN fig:ee-pi-plane
\begin{figure}[htbp]
\centering
\begin{tikzpicture}[fd,
  fa lab/.style={font=\footnotesize},
  fa contour/.style={draw=blue!65!black,line width=1.1pt,postaction={decorate,decoration={markings,
      mark=at position 0.13 with {\arrow{Stealth[length=6pt,width=5pt]}},
      mark=at position 0.5 with {\arrow{Stealth[length=6pt,width=5pt]}},
      mark=at position 0.87 with {\arrow{Stealth[length=6pt,width=5pt]}}}}},
  fa cutline/.style={draw=black,line width=1.0pt,decorate,decoration={zigzag,amplitude=1.7pt,segment length=4.6pt}},
  fa sm/.style={circle,draw=red!65!black,fill=white,line width=0.8pt,inner sep=0pt,minimum size=5.5pt}]
  \draw[lcaxis] (-4.0,0) -- (4.7,0) node[fa lab,anchor=west] {$\Real q^2$};
  \draw[lcaxis] (0,-3.15) -- (0,3.3) node[fa lab,left] {$\Imag q^2$};
  \draw[fa cutline] (0,0) -- (4.35,0);
  % contour of the contour formula: |q^2| = s, counter-clockwise from s+i0 to s-i0
  \draw[fa contour] (2.5:2.3) arc[start angle=2.5,end angle=357.5,radius=2.3];
  \node[fa lab,anchor=south west,inner sep=1.5pt] at (2.32,0.08) {$s+\ii0$};
  \node[fa lab,anchor=north west,inner sep=1.5pt] at (2.32,-0.08) {$s-\ii0$};
  \node[fa lab,text=blue!55!black,anchor=south east] at (135:2.38) {$|q^2|=s$};
  % smearing points s pm i Delta
  \draw[red!65!black,line width=0.5pt,dash pattern=on 1.5pt off 1.5pt] (2.3,-1.2) -- (2.3,1.2);
  \node[fa sm] at (2.3,1.2) {};
  \node[fa sm] at (2.3,-1.2) {};
  \node[fa lab,text=red!55!black,anchor=west] at (2.42,1.2) {$s+\ii\Delta$};
  \node[fa lab,text=red!55!black,anchor=west] at (2.42,-1.2) {$s-\ii\Delta$};
  % Euclidean point
  \fill (-1.3,0) circle[radius=2.2pt];
  \node[fa lab,anchor=north] at (-1.3,-0.08) {$-Q^2$};
  \node[fa lab,align=center,anchor=south] at (-1.3,0.18) {PT 可靠\\（欧氏区）};
  % cut and resonance region
  \draw[gray!60!black,line width=0.45pt] (3.9,0.1) -- (3.9,0.6);
  \node[fa lab,anchor=south,inner sep=1.5pt] at (3.9,0.6) {割线 $[0,\infty)$};
  \draw[gray!60!black,line width=0.5pt,-{Stealth[length=4pt,width=3.4pt]}] (0.82,-0.6) -- (0.36,-0.1);
  \node[fa lab,anchor=north west,align=left,inner sep=1pt] at (0.72,-0.56) {共振区\\（小 $s$）};
\end{tikzpicture}
\caption{流关联函数 $\Pi(q^2)$ 的解析结构：割线 $[0,\infty)$、欧氏点 $-Q^2$、围道公式 $R(s)=\frac1{2\pi}\int_{-\pi}^{\pi}\dd\theta\,D(s\ee^{\ii\theta})$ 的积分路径 $|q^2|=s$
（$D$ 的宗量 $Q^2=-q^2$ 走遍该圆周，在 $q^2$ 平面中即自 $s+\ii0$ 逆时针到 $s-\ii0$）以及抹平所用的点 $s\pm\ii\Delta$。在复 $q^2$ 处使用 OPE 是额外的假设（\cref{rem:ee-complex-ope}）。}
\label{fig:ee-pi-plane}
\end{figure}
%%% END fig:ee-pi-plane
